
\section{Conclusion}
\label{sec:conclusion}

This manuscript presented GibbsQ, a framework for entropy-regularized routing in heterogeneous parallel queues that addresses the tension between state-awareness, capacity-awareness, and formal stability guarantees. Our findings span three distinct layers of analysis:

At the \textbf{analytical level}, we formally demonstrated positive Harris recurrence for two routing laws via self-contained Foster--Lyapunov arguments (raw softmax and UAS), and established a parameter-conditional stability certificate for the SCUAS subclass. These proofs yield computable drift constants; notably, the UAS constants are concentration-parameter-independent, ensuring that stability does not degrade with increased entropy regularization.

At the \textbf{empirical level}, we showed that Calibrated UAS---the unrestricted three-parameter extension---achieves the lowest queue lengths among all evaluated closed-form policies, outperforming JSSQ by 8.8\% and baseline UAS by 12.6\%. We explicitly separate this empirical family from the certified SCUAS subclass, noting that follow-up auditing finds the benchmark-default setting to be uncertified by Theorem 6.7. However, the discovery of a numerically verified piecewise Lyapunov candidate for this default point provides concrete evidence for a future symbolic extension.

At the \textbf{learned-policy level}, we discovered a teacher-choice effect: RL agents initialized via behavior cloning from Calibrated UAS and fine-tuned via \reinforce{} surpass even their closed-form teacher under the Protocol B training distribution. This directional advantage is consistent across operating regimes and persists into near-critical loads up to $\rho = 0.98$.

The central practical takeaway is that the choice of analytical teacher during behavior cloning is the dominant design decision for neural queueing policies: the 1.32-unit gap between calibrated-teacher and UAS-teacher initialization exceeds all tested architectural modifications combined.

The most important limitation is that the default $\ngibbsq$ model (trained from a UAS teacher) does not outperform JSSQ or Calibrated UAS in the clean anchor benchmark; only the calibrated-teacher variant achieves this distinction. Additionally, the Calibrated UAS benchmark-default parameters remain uncertified by the SCUAS theorem, and all results assume Markovian (Poisson/exponential) dynamics.

The most promising theoretical direction is extending certification to competitive parameterizations via piecewise Lyapunov constructions (Appendix~\ref{sec:appendix_piecewise}), where computationally verified negative drift on exhaustive boundary shells provides strong evidence for a viable proof route. For practitioners, Calibrated UAS with $(\calbeta,\calgamma,\caloffset,\temp) = (0.85, 0.5, 0.5, 20)$ provides the strongest evaluated closed-form routing policy for heterogeneous parallel queues under the tested conditions.

For theorists, these results validate the tractability of soft-routing Lyapunov boundaries. For RL practitioners, they reinforce the critical value of analytically grounded inductive biases and high-quality teachers when optimizing sequential decisions in stochastic environments. Ultimately, this work demonstrates that analytically grounded routing policies provide a path from formal stability theory to effective neural policy training---a connection that neither analytical queueing theory nor unguided learned policies alone provides.
